% 原书分章 PDF 第 12—17 页：第 28.5 节，完整模拟设定、结果及五幅原数据图。
\section{模拟}

为评价靶向自助法的实际表现，我们从以下 $P_0$ 模拟数据：
\[
\begin{aligned}
 W_1&\sim U(-1,1),\\
 W_2&\sim U(-1,1),\\
 W^*&=W_2-W_1,\\
 A\mid W&\sim\mathrm{Bernoulli}(\operatorname{inv.logit}(-0.5W^*)),\\
 Y\mid A,W&\sim\mathrm{Bernoulli}
 (\operatorname{inv.logit}(A(1-0.5W^*+\sin(W^*)))).
\end{aligned}
\]
正值性，$(\bar g(P_0)(W)=P_0(A=1\mid W)>0$，得到满足：
$0.26<\bar g(P_0)(W)<0.74$。
在 $B=1000$ 次 Monte Carlo 模拟的每一次中，生成大小为 $n=1000$ 的样本。

模拟中，使用核回归估计 $\bar Q(P_0)(W)=E[Y\mid A=1,W]$，其带宽由十折交叉验证选择，即采用离散超级学习器。用设定正确的 logistic 回归估计 $\bar g(P_0)(W)$。每次模拟迭代中，先估计 $Q$，再拟合三种 TMLE：特定处理下均值的 TMLE、其渐近方差的 TMLE，以及同时估计特定处理下均值及其渐近方差的联合 TMLE。拟合这些 TMLE 后，采用五种不同方法，各重复生成 1000 个自助样本：非参数自助法、按簇的非参数自助法、基于超级学习器初始拟合的基于模型的自助法、从靶向于渐近方差的 TMLE 分布中抽样的靶向自助法，以及从联合靶向 TMLE 分布中抽样的联合靶向自助法。

对模拟数据集拟合的三种 TMLE，生成了不同的置信区间估计：第一种 TMLE 根据影响曲线构造 Wald 型区间，其余两种 TMLE 则直接估计方差。对于五种自助方法，我们均采用第 28.4.4 节介绍的全部四种方法估计区间。对所有估计的置信区间，评价其覆盖率和区间长度。我们还比较了超级学习器在完整样本和各类自助样本上的表现。为评价超级学习器在自助样本上的表现，比较了各种样本类型下选出的带宽，以及这些选择所对应的风险差比值。

\noindent\textbf{结果。}\enspace
如前所述，超级学习器在非参数自助样本上的行为与其在完整样本上不同，更接近回代估计量。在完整样本上，超级学习器，即交叉验证，通常选出接近预言选择器所选的带宽，后者使真实风险最小；而在非参数自助样本上，超级学习器最常选择可用带宽中的最小值，因而对数据过拟合。在包括靶向自助法在内的其他类型自助样本上，超级学习器的行为更接近其在完整样本上的行为，表明这些自助方法没有同样的问题。选择行为的这一差异，会影响所得超级学习器的风险差比值，见式 (28.1)。图 28.1 展示了这一点。其他自助样本在风险差比值方面，同样更接近完整样本。

\begin{figure}[htbp]
\centering
% 原书真实矢量数据及标记字形；中文只覆盖数据区外的标签。
\begin{tikzpicture}[x=0.51mm,y=-0.51mm]
\node[anchor=north west,inner sep=0pt] at (0,0) {\includegraphics[width=150.9600mm,height=98.9400mm]{figures/ch28_fig01_data.pdf}};
\node[font=\fontsize{9.5}{10.5}\selectfont,inner sep=0pt,anchor=east] at (88.24687,155.37414) {联合靶向自助法};
\node[font=\fontsize{9.5}{10.5}\selectfont,inner sep=0pt,anchor=east] at (88.24831,128.15646) {靶向自助法};
\node[font=\fontsize{9.5}{10.5}\selectfont,inner sep=0pt,anchor=east] at (88.24117,100.94662) {超级学习自助法};
\node[font=\fontsize{9.5}{10.5}\selectfont,inner sep=0pt,anchor=east] at (88.24329,73.72897) {非参数自助法 + CVID};
\node[font=\fontsize{9.5}{10.5}\selectfont,inner sep=0pt,anchor=east] at (88.24332,46.51773) {非参数自助法};
\node[font=\fontsize{9.5}{10.5}\selectfont,inner sep=0pt,anchor=east] at (88.24402,19.30008) {完整样本};
\node[font=\fontsize{9.5}{10.5}\selectfont,inner sep=0pt,anchor=center] at (193.41423,187.65699) {风险差比值};
\node[font=\fontsize{9.5}{10.5}\selectfont,inner sep=0pt,anchor=center,rotate=90] at (6.86757,86.71398) {样本类型};
\end{tikzpicture}

\caption{超级学习器在不同样本类型上的风险差比值（定义见式 (28.1)）。}
\label{fig:ch28-risk-ratio}
\end{figure}

图 28.2、28.3 和 28.4 展示了超级学习器的表现如何影响所得 TMLE 估计的置信区间表现。一般而言，在非参数自助样本上用于 TMLE 的过拟合超级学习器，比在完整样本上用于 TMLE 的良好拟合超级学习器具有更大的变异性。因此，非参数自助法置信区间不必要地变长，并产生过度覆盖。在 $n=1000$ 时，这种过度覆盖对长度的影响不大：由非参数自助法估计的 Wald 区间，平均仅比标准的基于影响曲线的置信区间长 $4\%$。在较小样本量下，影响更严重：当 $n=250$ 时，非参数自助法区间比影响曲线区间长 $21\%$；当 $n=100$ 时，则长 $39\%$。长度的这种显著增加，会不利于非参数自助法置信区间的统计功效。在我们的模拟中，超级学习器可选择的带宽集合不随样本量变化。我们预期，如果还可以选择更小的带宽，那么超级学习器在非参数自助样本上就会选择它们，从而使过拟合在较大样本量下的影响进一步增加。

\clearpage
\begin{landscape}
\begin{figure}[p]
\centering
% 原书真实矢量数据及标记字形；中文只覆盖数据区外的标签。
\begin{tikzpicture}[x=0.56mm,y=-0.56mm]
\node[anchor=north west,inner sep=0pt] at (0,0) {\includegraphics[width=188.1600mm,height=140.0000mm]{figures/ch28_fig02_data.pdf}};
\node[font=\fontsize{9.5}{10.5}\selectfont,inner sep=0pt,anchor=center] at (130.32292,119.25395) {双侧覆盖率};
\node[font=\fontsize{9.5}{10.5}\selectfont,inner sep=0pt,anchor=center] at (266.08897,119.25395) {区间长度};
\node[font=\fontsize{9.5}{10.5}\selectfont,inner sep=0pt,anchor=center] at (130.31863,7.85616) {下侧覆盖率};
\node[font=\fontsize{9.5}{10.5}\selectfont,inner sep=0pt,anchor=center] at (266.09593,7.85616) {上侧覆盖率};
\node[font=\fontsize{8.5}{9.5}\selectfont,inner sep=0pt,anchor=east] at (61.12234,100.35838) {影响曲线};
\node[font=\fontsize{8.5}{9.5}\selectfont,inner sep=0pt,anchor=east] at (60.00256,89.19933) {靶向方差};
\node[font=\fontsize{8.5}{9.5}\selectfont,inner sep=0pt,anchor=east] at (59.99239,78.04240) {联合靶向方差};
\node[font=\fontsize{8.5}{9.5}\selectfont,inner sep=0pt,anchor=east] at (59.98968,66.88927) {非参数自助法};
\node[font=\fontsize{8.5}{9.5}\selectfont,inner sep=0pt,anchor=east] at (59.99076,55.73073) {非参数自助法 + CVID};
\node[font=\fontsize{8.5}{9.5}\selectfont,inner sep=0pt,anchor=east] at (59.99398,44.57167) {超级学习自助法};
\node[font=\fontsize{8.5}{9.5}\selectfont,inner sep=0pt,anchor=east] at (59.99723,33.42067) {靶向自助法};
\node[font=\fontsize{8.5}{9.5}\selectfont,inner sep=0pt,anchor=east] at (59.99830,22.26161) {联合靶向自助法};
\node[font=\fontsize{8.5}{9.5}\selectfont,inner sep=0pt,anchor=east] at (61.12312,211.75697) {影响曲线};
\node[font=\fontsize{8.5}{9.5}\selectfont,inner sep=0pt,anchor=east] at (60.00335,200.59792) {靶向方差};
\node[font=\fontsize{8.5}{9.5}\selectfont,inner sep=0pt,anchor=east] at (59.99318,189.44582) {联合靶向方差};
\node[font=\fontsize{8.5}{9.5}\selectfont,inner sep=0pt,anchor=east] at (59.99046,178.28676) {非参数自助法};
\node[font=\fontsize{8.5}{9.5}\selectfont,inner sep=0pt,anchor=east] at (59.99155,167.12877) {非参数自助法 + CVID};
\node[font=\fontsize{8.5}{9.5}\selectfont,inner sep=0pt,anchor=east] at (59.99477,155.97670) {超级学习自助法};
\node[font=\fontsize{8.5}{9.5}\selectfont,inner sep=0pt,anchor=east] at (59.99802,144.81764) {靶向自助法};
\node[font=\fontsize{8.5}{9.5}\selectfont,inner sep=0pt,anchor=east] at (59.99909,133.65965) {联合靶向自助法};
\node[font=\fontsize{9.5}{10.5}\selectfont,inner sep=0pt,anchor=center] at (135.92426,245.58475) {区间类型};
\node[font=\fontsize{8.5}{9.5}\selectfont,inner sep=0pt,anchor=center] at (173.93143,245.48196) {BCa};
\node[font=\fontsize{8.5}{9.5}\selectfont,inner sep=0pt,anchor=center] at (205.54599,245.48196) {bootstrap-$t$};
\node[font=\fontsize{8.5}{9.5}\selectfont,inner sep=0pt,anchor=center] at (243.57878,245.48196) {百分位};
\node[font=\fontsize{8.5}{9.5}\selectfont,inner sep=0pt,anchor=center] at (273.59612,245.48196) {Wald};
\end{tikzpicture}

\caption{$n=1000$ 时置信区间的覆盖率与长度。}
\label{fig:ch28-ci-1000}
\end{figure}
\end{landscape}

\begin{landscape}
\begin{figure}[p]
\centering
% 原书真实矢量数据及标记字形；中文只覆盖数据区外的标签。
\begin{tikzpicture}[x=0.645mm,y=-0.645mm]
\node[anchor=north west,inner sep=0pt] at (0,0) {\includegraphics[width=188.3400mm,height=140.6100mm]{figures/ch28_fig03_data.pdf}};
\node[font=\fontsize{9.5}{10.5}\selectfont,inner sep=0pt,anchor=center] at (112.38167,103.72179) {双侧覆盖率};
\node[font=\fontsize{9.5}{10.5}\selectfont,inner sep=0pt,anchor=center] at (230.22197,103.72179) {区间长度};
\node[font=\fontsize{9.5}{10.5}\selectfont,inner sep=0pt,anchor=center] at (112.37670,7.03161) {下侧覆盖率};
\node[font=\fontsize{9.5}{10.5}\selectfont,inner sep=0pt,anchor=center] at (230.22984,7.03161) {上侧覆盖率};
\node[font=\fontsize{8.5}{9.5}\selectfont,inner sep=0pt,anchor=east] at (52.31801,87.32171) {影响曲线};
\node[font=\fontsize{8.5}{9.5}\selectfont,inner sep=0pt,anchor=east] at (52.32313,77.63592) {靶向方差};
\node[font=\fontsize{8.5}{9.5}\selectfont,inner sep=0pt,anchor=east] at (52.32080,67.95153) {联合靶向方差};
\node[font=\fontsize{8.5}{9.5}\selectfont,inner sep=0pt,anchor=east] at (52.31429,58.27087) {非参数自助法};
\node[font=\fontsize{8.5}{9.5}\selectfont,inner sep=0pt,anchor=east] at (52.31476,48.58599) {非参数自助法 + CVID};
\node[font=\fontsize{8.5}{9.5}\selectfont,inner sep=0pt,anchor=east] at (52.31802,38.90069) {超级学习自助法};
\node[font=\fontsize{8.5}{9.5}\selectfont,inner sep=0pt,anchor=east] at (52.32221,29.22100) {靶向自助法};
\node[font=\fontsize{8.5}{9.5}\selectfont,inner sep=0pt,anchor=east] at (52.32547,19.53570) {联合靶向自助法};
\node[font=\fontsize{8.5}{9.5}\selectfont,inner sep=0pt,anchor=east] at (52.31813,184.01278) {影响曲线};
\node[font=\fontsize{8.5}{9.5}\selectfont,inner sep=0pt,anchor=east] at (52.32326,174.32793) {靶向方差};
\node[font=\fontsize{8.5}{9.5}\selectfont,inner sep=0pt,anchor=east] at (52.32093,164.64776) {联合靶向方差};
\node[font=\fontsize{8.5}{9.5}\selectfont,inner sep=0pt,anchor=east] at (52.31441,154.96291) {非参数自助法};
\node[font=\fontsize{8.5}{9.5}\selectfont,inner sep=0pt,anchor=east] at (52.31489,145.27712) {非参数自助法 + CVID};
\node[font=\fontsize{8.5}{9.5}\selectfont,inner sep=0pt,anchor=east] at (52.31814,135.59695) {超级学习自助法};
\node[font=\fontsize{8.5}{9.5}\selectfont,inner sep=0pt,anchor=east] at (52.32233,125.91210) {靶向自助法};
\node[font=\fontsize{8.5}{9.5}\selectfont,inner sep=0pt,anchor=east] at (52.32559,116.22726) {联合靶向自助法};
\node[font=\fontsize{9.5}{10.5}\selectfont,inner sep=0pt,anchor=center] at (117.24309,213.37561) {区间类型};
\node[font=\fontsize{8.5}{9.5}\selectfont,inner sep=0pt,anchor=center] at (150.23354,213.28580) {BCa};
\node[font=\fontsize{8.5}{9.5}\selectfont,inner sep=0pt,anchor=center] at (177.67186,213.28580) {bootstrap-$t$};
\node[font=\fontsize{8.5}{9.5}\selectfont,inner sep=0pt,anchor=center] at (210.68391,213.28580) {百分位};
\node[font=\fontsize{8.5}{9.5}\selectfont,inner sep=0pt,anchor=center] at (236.74223,213.28580) {Wald};
\end{tikzpicture}

\caption{$n=250$ 时置信区间的覆盖率与长度。}
\label{fig:ch28-ci-250}
\end{figure}
\end{landscape}

\begin{landscape}
\begin{figure}[p]
\centering
% 原书真实矢量数据及标记字形；中文只覆盖数据区外的标签。
\begin{tikzpicture}[x=0.645mm,y=-0.645mm]
\node[anchor=north west,inner sep=0pt] at (0,0) {\includegraphics[width=188.3400mm,height=140.6100mm]{figures/ch28_fig04_data.pdf}};
\node[font=\fontsize{9.5}{10.5}\selectfont,inner sep=0pt,anchor=center] at (113.39987,104.00481) {双侧覆盖率};
\node[font=\fontsize{9.5}{10.5}\selectfont,inner sep=0pt,anchor=center] at (230.82796,104.00481) {区间长度};
\node[font=\fontsize{9.5}{10.5}\selectfont,inner sep=0pt,anchor=center] at (113.15598,7.31459) {下侧覆盖率};
\node[font=\fontsize{9.5}{10.5}\selectfont,inner sep=0pt,anchor=center] at (230.94531,7.31459) {上侧覆盖率};
\node[font=\fontsize{8.5}{9.5}\selectfont,inner sep=0pt,anchor=east] at (52.92735,87.60466) {影响曲线};
\node[font=\fontsize{8.5}{9.5}\selectfont,inner sep=0pt,anchor=east] at (52.92735,77.91795) {靶向方差};
\node[font=\fontsize{8.5}{9.5}\selectfont,inner sep=0pt,anchor=east] at (52.92875,68.23125) {联合靶向方差};
\node[font=\fontsize{8.5}{9.5}\selectfont,inner sep=0pt,anchor=east] at (52.92736,58.54967) {非参数自助法};
\node[font=\fontsize{8.5}{9.5}\selectfont,inner sep=0pt,anchor=east] at (52.92737,48.86201) {非参数自助法 + CVID};
\node[font=\fontsize{8.5}{9.5}\selectfont,inner sep=0pt,anchor=east] at (52.92735,39.17531) {超级学习自助法};
\node[font=\fontsize{8.5}{9.5}\selectfont,inner sep=0pt,anchor=east] at (52.92688,29.49373) {靶向自助法};
\node[font=\fontsize{8.5}{9.5}\selectfont,inner sep=0pt,anchor=east] at (52.92781,19.80702) {联合靶向自助法};
\node[font=\fontsize{8.5}{9.5}\selectfont,inner sep=0pt,anchor=east] at (52.92461,184.29582) {影响曲线};
\node[font=\fontsize{8.5}{9.5}\selectfont,inner sep=0pt,anchor=east] at (52.92460,174.60863) {靶向方差};
\node[font=\fontsize{8.5}{9.5}\selectfont,inner sep=0pt,anchor=east] at (52.92600,164.92656) {联合靶向方差};
\node[font=\fontsize{8.5}{9.5}\selectfont,inner sep=0pt,anchor=east] at (52.92461,155.24080) {非参数自助法};
\node[font=\fontsize{8.5}{9.5}\selectfont,inner sep=0pt,anchor=east] at (52.92462,145.55363) {非参数自助法 + CVID};
\node[font=\fontsize{8.5}{9.5}\selectfont,inner sep=0pt,anchor=east] at (52.92461,135.87065) {超级学习自助法};
\node[font=\fontsize{8.5}{9.5}\selectfont,inner sep=0pt,anchor=east] at (52.92460,126.18486) {靶向自助法};
\node[font=\fontsize{8.5}{9.5}\selectfont,inner sep=0pt,anchor=east] at (52.92553,116.49767) {联合靶向自助法};
\node[font=\fontsize{9.5}{10.5}\selectfont,inner sep=0pt,anchor=center] at (118.17627,213.65556) {区间类型};
\node[font=\fontsize{8.5}{9.5}\selectfont,inner sep=0pt,anchor=center] at (150.83955,213.56885) {BCa};
\node[font=\fontsize{8.5}{9.5}\selectfont,inner sep=0pt,anchor=center] at (178.30001,213.56885) {bootstrap-$t$};
\node[font=\fontsize{8.5}{9.5}\selectfont,inner sep=0pt,anchor=center] at (211.28687,213.56885) {百分位};
\node[font=\fontsize{8.5}{9.5}\selectfont,inner sep=0pt,anchor=center] at (237.37878,213.56885) {Wald};
\end{tikzpicture}

\caption{$n=100$ 时置信区间的覆盖率与长度。}
\label{fig:ch28-ci-100}
\end{figure}
\end{landscape}

这些图也表明，自助法同时靶向于感兴趣参数及其渐近方差十分重要。对于 Wald 以外的区间类型，仅靶向于方差的靶向自助法区间，其覆盖率很差。这是因为这些区间的中心，并非完整数据集得到的特定处理下均值估计，而是自助区间所得到的估计的平均值。对于靶向自助样本，这些估计存在偏差，因为靶向自助法只靶向于方差，没有靶向于真正感兴趣的参数。

图 28.4 表明，在较小样本量下，渐近 Wald 区间的覆盖率低于名义水平，所有方法都至少少覆盖 $2.5\%$。在这样的小样本中，自助法最有可能改进渐近置信区间。样本量较大时，二阶项变得相对不重要。然而，即使在这样的小样本量下，本模拟中的渐近区间也只是略微偏于反保守。这可能是因为，即使 $n=100$，模拟的抽样分布也已十分接近正态。

只关注联合靶向自助法，就可以比较不同类型自助法置信区间的表现，见图 28.5。样本量不大时，bootstrap-$t$ 区间过度覆盖，而且比其他区间类型更长。其他自助方法生成的区间长度相近。在这三种方法中，$BC_a$ 在所考察的样本量范围内，其覆盖率最接近名义水平。因此，建议今后将这种区间类型与联合靶向自助法结合使用。

\begin{landscape}
\begin{figure}[p]
\centering
% 原书真实矢量数据及标记字形；中文只覆盖数据区外的标签。
\begin{tikzpicture}[x=0.635mm,y=-0.635mm]
\node[anchor=north west,inner sep=0pt] at (0,0) {\includegraphics[width=185.4200mm,height=140.9700mm]{figures/ch28_fig05_data.pdf}};
\node[font=\fontsize{9.5}{10.5}\selectfont,inner sep=0pt,anchor=center] at (67.22526,110.77368) {双侧覆盖率};
\node[font=\fontsize{9.5}{10.5}\selectfont,inner sep=0pt,anchor=center] at (188.95396,110.77368) {区间长度};
\node[font=\fontsize{9.5}{10.5}\selectfont,inner sep=0pt,anchor=center] at (67.21775,6.67758) {下侧覆盖率};
\node[font=\fontsize{9.5}{10.5}\selectfont,inner sep=0pt,anchor=center] at (188.95332,6.67758) {上侧覆盖率};
\node[font=\fontsize{9.5}{10.5}\selectfont,inner sep=0pt,anchor=center] at (267.54095,90.41000) {区间类型};
\node[font=\fontsize{8.5}{9.5}\selectfont,inner sep=0pt,anchor=west] at (261.47000,99.08279) {BCa};
\node[font=\fontsize{8.5}{9.5}\selectfont,inner sep=0pt,anchor=west] at (261.47000,108.03680) {bootstrap-$t$};
\node[font=\fontsize{8.5}{9.5}\selectfont,inner sep=0pt,anchor=west] at (261.47000,116.99176) {百分位};
\node[font=\fontsize{8.5}{9.5}\selectfont,inner sep=0pt,anchor=west] at (261.47000,125.94580) {Wald};
\end{tikzpicture}

\caption{联合靶向自助法区间的表现比较。}
\label{fig:ch28-joint-comparison}
\end{figure}
\end{landscape}
